\documentclass[aspectratio=169,11pt]{beamer}

\usetheme{Madrid}
\usecolortheme{default}
\setbeamertemplate{navigation symbols}{}

\usepackage[T1]{fontenc}
\usepackage[utf8]{inputenc}
\usepackage{lmodern}
\usepackage{amsmath, amssymb}
\usepackage{booktabs}
\usepackage{pgfplots}
\usepackage{pgfplotstable}
\usepackage{siunitx}
\usepackage{hyperref}

\pgfplotsset{compat=1.18}
\sisetup{group-separator={,}, round-mode=places, round-precision=2}

\newcommand{\CourseName}{Course Name Placeholder}
\newcommand{\InstructorName}{Instructor Placeholder}
\newcommand{\StudentName}{Student Name Placeholder}
\newcommand{\StudentID}{Student ID Placeholder}

\title[Digital Twin SimOpt]{Project 2: Digital Twin-Enabled Simulation-Based Optimization}
\subtitle{Continuous Calibration and Simulation-Based Decision Selection in Plant Simulation}
\author{\StudentName \\ \small \StudentID}
\institute{\CourseName \\ \small Instructor: \InstructorName}
\date{\today}

\begin{document}

\begin{frame}
    \titlepage
\end{frame}

\begin{frame}{Presentation Roadmap}
    \tableofcontents
\end{frame}

\section{Problem}

\begin{frame}{Problem Definition}
    \textbf{Goal:} use the synchronized digital twin from Project 1 to optimize decisions online instead of only monitoring the system.

    \vspace{0.4cm}
    \textbf{Objective}
    \[
    \max_X \; J_t(X;\theta_t)=\mathbb{E}[\text{NetRevenue}\mid X,\theta_t]
    \]

    \textbf{Decision variables}
    \[
    X=(Q,R,W,M1,M2)
    \]
    with
    \[
    Q,R \in \{50,\dots,100\},\;
    W \in \{10,\dots,15\},\;
    M1,M2 \in \{1,2,3\}
    \]

    \textbf{Key idea}
    \begin{itemize}
        \item Project 1 estimates the latent arrival parameter $\theta$ online.
        \item Project 2 uses that synchronized state to select high-revenue decisions.
    \end{itemize}
\end{frame}

\begin{frame}{Bridge from Project 1 to Project 2}
    \textbf{Project 1 output:} online posterior over the customer-arrival parameter $\theta$.

    \vspace{0.4cm}
    \textbf{Project 2 loop}
    \begin{enumerate}
        \item observe new \((X,y)\) data from the physical/proxy system,
        \item recalibrate the digital twin online,
        \item use the synchronized twin to score candidate decisions,
        \item deploy new candidates and validate them on the physical system.
    \end{enumerate}

    \vspace{0.4cm}
    \textbf{Why this matters}
    \begin{itemize}
        \item a static nominal policy may become suboptimal when $\theta$ changes,
        \item a synchronized twin can adapt optimization decisions to changing conditions.
    \end{itemize}
\end{frame}

\section{Method}

\begin{frame}{Method Overview}
    \textbf{Digital twin components}
    \begin{itemize}
        \item offline neural emulator trained on aggregated Plant Simulation data,
        \item online restart BOCPD + particle filter for continuous calibration.
    \end{itemize}

    \vspace{0.3cm}
    \textbf{Selection policy}
    \[
    \text{score}(X)=\mu(X)+\beta\sigma(X), \qquad \beta=0.3
    \]
    where $\mu(X)$ and $\sigma(X)$ come from the synchronized twin.

    \vspace{0.3cm}
    \textbf{Per batch}
    \begin{itemize}
        \item select top-5 candidate points,
        \item compare with fixed baseline top-5 from $\theta=11.5$,
        \item evaluate 10 points in Plant Simulation,
        \item feed all 10 observations back into calibration.
    \end{itemize}
\end{frame}

\begin{frame}{Experimental Design}
    \textbf{Nonstationary environment}
    \begin{itemize}
        \item true $\theta$ switches among 4 platforms:
        \[
        \{3,\;11.5,\;20,\;30\}
        \]
        \item each platform lasts 20 batches,
        \item within each platform:
        \[
        \theta_t \sim \mathcal{N}(\theta_{\text{platform}}, 0.8^2), \quad \theta_t > 0
        \]
    \end{itemize}

    \vspace{0.3cm}
    \textbf{Why this schedule?}
    \begin{itemize}
        \item platform jumps create clearer regime changes than a smooth sinusoidal drift,
        \item easier to test whether synchronization improves decision quality.
    \end{itemize}

    \vspace{0.3cm}
    \textbf{Run configuration}
    \begin{itemize}
        \item 80 batches,
        \item 512 particles,
        \item 5 candidate points + 5 baseline points per batch,
        \item backend: Plant Simulation.
    \end{itemize}
\end{frame}

\section{Results}

\begin{frame}{Theta Tracking}
    \begin{center}
    \begin{tikzpicture}
    \begin{axis}[
        width=0.95\textwidth,
        height=0.58\textheight,
        xlabel={Batch index},
        ylabel={$\theta$ (minutes)},
        grid=both,
        legend style={at={(0.5,1.02)}, anchor=south, legend columns=2},
    ]
    \addplot[blue, thick] table[x=batch_idx, y=theta_true, col sep=comma] {project2_batch_summary.csv};
    \addlegendentry{True $\theta$}
    \addplot[red, thick] table[x=batch_idx, y=theta_est, col sep=comma] {project2_batch_summary.csv};
    \addlegendentry{Estimated $\theta$}
    \end{axis}
    \end{tikzpicture}
    \end{center}

    \textbf{Takeaway:}
    the calibrator tracks the platform changes reasonably well, but visible lag appears after major jumps.
\end{frame}

\begin{frame}{Main Quantitative Results}
    \begin{table}
    \centering
    \begin{tabular}{p{0.62\linewidth}r}
    \toprule
    Metric & Value \\
    \midrule
    Final cumulative mean improvement over baseline & 9,251,549.6 \\
    Final cumulative max improvement over baseline & 8,326,897.0 \\
    Average absolute theta tracking error & 1.31 \\
    Median absolute theta tracking error & 0.96 \\
    Restart count & 3 \\
    Candidate mean better than baseline mean & 27 / 80 \\
    Candidate max better than baseline max & 19 / 80 \\
    \bottomrule
    \end{tabular}
    \end{table}

    \textbf{Interpretation}
    \begin{itemize}
        \item the synchronized policy wins in cumulative performance,
        \item but it does not dominate the baseline in every single batch,
        \item synchronization helps most over the full horizon rather than instantaneously.
    \end{itemize}
\end{frame}

\begin{frame}{Optimization Progress}
    \begin{center}
    \begin{tikzpicture}
    \begin{axis}[
        width=0.95\textwidth,
        height=0.58\textheight,
        xlabel={Batch index},
        ylabel={Cumulative gap},
        grid=both,
        legend style={at={(0.5,1.02)}, anchor=south, legend columns=2},
    ]
    \addplot[teal, thick] table[x=batch_idx, y=cum_gap_mean, col sep=comma] {project2_batch_summary.csv};
    \addlegendentry{Candidate mean - baseline mean}
    \addplot[orange, thick] table[x=batch_idx, y=cum_gap_best, col sep=comma] {project2_batch_summary.csv};
    \addlegendentry{Candidate max - baseline max}
    \end{axis}
    \end{tikzpicture}
    \end{center}

    \textbf{Takeaway:}
    despite temporary losses after some jumps, the synchronized policy accumulates a strong positive advantage over time.
\end{frame}

\begin{frame}{Reality Gap and Validation}
    \textbf{Validation step}
    \begin{itemize}
        \item decisions selected by the twin are directly evaluated in Plant Simulation,
        \item this closes the loop between prediction and physical/proxy validation.
    \end{itemize}

    \vspace{0.4cm}
    \textbf{Reality gap}
    \[
    \text{RealityGap}(X_t)=y_t^{\text{actual}}-y_t^{\text{pred}}
    \]

    \begin{itemize}
        \item average reality gap (mean reward): \(-45{,}827.57\)
        \item average reality gap (max reward): \(-42{,}831.38\)
    \end{itemize}

    \textbf{Interpretation}
    \begin{itemize}
        \item the twin is somewhat optimistic on selected points,
        \item but cumulative gains remain positive even after physical validation.
    \end{itemize}
\end{frame}

\section{Discussion}

\begin{frame}{Discussion}
    \textbf{What worked}
    \begin{itemize}
        \item the digital twin tracked regime switches with moderate error,
        \item synchronization produced positive cumulative improvement over a nominal baseline,
        \item the framework successfully closed the calibration--optimization--validation loop.
    \end{itemize}

    \vspace{0.3cm}
    \textbf{Main limitations}
    \begin{itemize}
        \item candidate search is pool-based rather than globally optimal,
        \item selected points still show an optimism-driven reality gap,
        \item baseline remains competitive in many single batches.
    \end{itemize}

    \vspace{0.3cm}
    \textbf{Practical implication}
    \begin{itemize}
        \item synchronization is most valuable in rolling-horizon planning under nonstationarity.
    \end{itemize}
\end{frame}

\begin{frame}{Conclusion}
    \textbf{Conclusion}
    \begin{itemize}
        \item Project 1 provided a synchronized digital twin.
        \item Project 2 used that synchronized twin as a decision-support engine.
        \item Under platform-switching demand conditions, the proposed policy outperformed a fixed baseline in cumulative reward.
    \end{itemize}

    \vspace{0.4cm}
    \textbf{Main message}
    \begin{block}{}
    A synchronized digital twin can support simulation-based optimization under changing operating conditions, not just monitoring.
    \end{block}

    \vspace{0.3cm}
    \textbf{Future work}
    \begin{itemize}
        \item improve ranking accuracy of the twin,
        \item compare posterior-UCB with direct $\theta_{\text{est}}$-conditioned exploitation,
        \item refine candidate search over the full discrete space.
    \end{itemize}
\end{frame}

\begin{frame}{Thank You}
    \centering
    {\Large Thank you!}

    \vspace{0.6cm}
    Questions and discussion.
\end{frame}

\end{document}
